% GENERATED FILE. Edit canonical lessons, then run npm run generate:books.
% canonical-source-hash: fnv1a64:fb57890eb5105046
% canonical-lessons: JA-C330-chanto, JA-C330-shikkari, JA-C330-minarau, JA-C330-soroeru, JA-C330-aratameru

\chapter{Doing Things Properly}
\label{ch:ja-doing-things-properly}
\hlchaptermodality{330}

\begin{chapteropening}
\textbf{By the end of this chapter:} \emph{I can tell someone to do a thing properly and say why.}

Complete the last lesson of chapter 330: I can tell someone to do a thing properly and say why.
\end{chapteropening}

\section[\texorpdfstring{\ja{ちゃんと} (chanto)}{ちゃんと (chanto)}]{\ja{ちゃんと} (chanto) — as one should; properly}

\label{lesson:JA-C330-chanto}



\begin{quote}
Before the new one: say the Japanese for one's turn, then the Japanese for a cigarette end.
\end{quote}

\subsection*{You'll want to know: \ja{ちゃんと}}
\textbf{\ja{ちゃんと}} — \emph{chanto} — \textquotedblleft{}as one should; properly\textquotedblright{}.

\begin{quote}
\textbf{\ja{びょうきに} \ja{なるから、ちゃんと} \ja{てを} \ja{あらって} \ja{ください。}} — \emph{byōki ni naru kara, chanto te o aratte kudasai} — \textquotedblleft{}You will get ill, so wash your hands properly.\textquotedblright{}
\end{quote}

\begin{cousinweb}
\textbf{\ja{ちゃんと}} (\emph{chanto}) is a native mimetic word, one of the Japanese words that picture a manner or a feeling in sound; JMdict marks it as mimetic. It is inherited, not borrowed, and its older origin is not settled (JMdict).
\end{cousinweb}

\begin{grammarlens}[title={What you've built}]
One more word for giving reasons and opinions.
\end{grammarlens}

\subsection*{Your turn}
\begin{itemize}
\raggedright
  \item \emph{Say it:} \emph{chanto}
  \item \emph{Say it:} \emph{chanto}, once more
  \item \emph{Say it:} \emph{byōki ni naru kara, chanto te o aratte kudasai} — You will get ill, so wash your hands properly.
  \item \emph{Recall:} say the Japanese for one's turn, then the Japanese for a cigarette end, then say \emph{chanto} again
\end{itemize}

\begin{tcolorbox}[breakable,colback=teal!4,colframe=teal!35!black,title={Before you move on}]
What does \ja{ちゃんと} mean? (\textquotedblleft{}As one should; properly\textquotedblright{}.) Say it once more.
\end{tcolorbox}
\section[\texorpdfstring{\ja{しっかり} (shikkari)}{しっかり (shikkari)}]{\ja{しっかり} (shikkari) — firmly; steadily}

\label{lesson:JA-C330-shikkari}



\begin{quote}
Before the new one: say the Japanese for a cigarette end, then the Japanese for as one should.
\end{quote}

\subsection*{You'll want to know: \ja{しっかり}}
\textbf{\ja{しっかり}} — \emph{shikkari} — \textquotedblleft{}firmly; steadily\textquotedblright{}.

\begin{quote}
\textbf{\ja{かぜが} \ja{つよいから、しっかり} \ja{もって} \ja{ください。}} — \emph{kaze ga tsuyoi kara, shikkari motte kudasai} — \textquotedblleft{}The wind is strong, so hold on firmly.\textquotedblright{}
\end{quote}

\begin{cousinweb}
\textbf{\ja{しっかり}} (\emph{shikkari}) is a native mimetic word, one of the Japanese words that picture a manner or a feeling in sound; JMdict marks it as mimetic. It is inherited, not borrowed, and its older origin is not settled (JMdict).
\end{cousinweb}

\begin{grammarlens}[title={What you've built}]
One more word for giving reasons and opinions.
\end{grammarlens}

\subsection*{Your turn}
\begin{itemize}
\raggedright
  \item \emph{Say it:} \emph{shikkari}
  \item \emph{Say it:} \emph{shikkari}, once more
  \item \emph{Say it:} \emph{kaze ga tsuyoi kara, shikkari motte kudasai} — The wind is strong, so hold on firmly.
  \item \emph{Recall:} say the Japanese for a cigarette end, then the Japanese for as one should, then say \emph{shikkari} again
\end{itemize}

\begin{tcolorbox}[breakable,colback=teal!4,colframe=teal!35!black,title={Before you move on}]
What does \ja{しっかり} mean? (\textquotedblleft{}Firmly; steadily\textquotedblright{}.) Say it once more.
\end{tcolorbox}
\section[\texorpdfstring{\ja{みならう} (minarau)}{みならう (minarau)}]{\ja{みならう} (minarau) — to follow an example; to model oneself on}

\label{lesson:JA-C330-minarau}



\begin{quote}
Before the new one: say the Japanese for as one should, then the Japanese for firmly.
\end{quote}

\subsection*{You'll want to know: \ja{みならう}}
\textbf{\ja{みならう}} — \emph{minarau} — \textquotedblleft{}to follow an example; to model oneself on\textquotedblright{}.

\begin{quote}
\textbf{\ja{あには} \ja{よく} \ja{べんきょうするから、みならった} \ja{ほうが} \ja{いいです。}} — \emph{ani wa yoku benkyō suru kara, minaratta hō ga ii desu} — \textquotedblleft{}Your brother studies hard, so you should follow his example.\textquotedblright{}
\end{quote}

Polite \textbf{\ja{みならいます}}, \emph{minaraimasu}; a godan verb. Past (polite): \textbf{\ja{おとうとは} \ja{あにを} \ja{みならいました。}}, \emph{otōto wa ani o minaraimashita}.

\begin{cousinweb}
\textbf{\ja{みならう}} (\emph{minarau}) is a native Japanese word (wago), inherited rather than borrowed from Chinese: it is built from \emph{mi}, \textquotedblleft{}seeing (miru)\textquotedblright{} and \emph{narau}, \textquotedblleft{}to learn\textquotedblright{}, and its kanji were attached later as kun'yomi (JMdict).
\end{cousinweb}

\begin{grammarlens}[title={What you've built}]
One more word for giving reasons and opinions.
\end{grammarlens}

\subsection*{Your turn}
\begin{itemize}
\raggedright
  \item \emph{Say it:} \emph{minarau}
  \item \emph{Say it:} \emph{minarau}, once more
  \item \emph{Say it:} \emph{ani wa yoku benkyō suru kara, minaratta hō ga ii desu} — Your brother studies hard, so you should follow his example.
  \item \emph{Recall:} say the Japanese for as one should, then the Japanese for firmly, then say \emph{minarau} again
\end{itemize}

\begin{tcolorbox}[breakable,colback=teal!4,colframe=teal!35!black,title={Before you move on}]
What does \ja{みならう} mean? (\textquotedblleft{}To follow an example; to model oneself on\textquotedblright{}.) Say it once more.
\end{tcolorbox}
\section[\texorpdfstring{\ja{そろえる} (soroeru)}{そろえる (soroeru)}]{\ja{そろえる} (soroeru) — to line up neatly; to put in order}

\label{lesson:JA-C330-soroeru}



\begin{quote}
Before the new one: say the Japanese for firmly, then the Japanese for to follow an example.
\end{quote}

\subsection*{You'll want to know: \ja{そろえる}}
\textbf{\ja{そろえる}} — \emph{soroeru} — \textquotedblleft{}to line up neatly; to put in order\textquotedblright{}.

\begin{quote}
\textbf{\ja{げんかんで} \ja{くつを} \ja{そろえた} \ja{ほうが} \ja{いいです。}} — \emph{genkan de kutsu o soroeta hō ga ii desu} — \textquotedblleft{}You should line your shoes up neatly in the entrance.\textquotedblright{}
\end{quote}

Polite \textbf{\ja{そろえます}}, \emph{soroemasu}; an ichidan verb. Past (polite): \textbf{\ja{こどもが} \ja{くつを} \ja{そろえました。}}, \emph{kodomo ga kutsu o soroemashita}.

\begin{cousinweb}
\textbf{\ja{そろえる}} (\emph{soroeru}) is a native Japanese word (wago), inherited rather than borrowed from Chinese, and its older origin is not settled. Its kanji was attached to it later as a kun'yomi, a native reading given to a Chinese character (JMdict, KANJIDIC2).
\end{cousinweb}

\begin{grammarlens}[title={What you've built}]
One more word for giving reasons and opinions.
\end{grammarlens}

\subsection*{Your turn}
\begin{itemize}
\raggedright
  \item \emph{Say it:} \emph{soroeru}
  \item \emph{Say it:} \emph{soroeru}, once more
  \item \emph{Say it:} \emph{genkan de kutsu o soroeta hō ga ii desu} — You should line your shoes up neatly in the entrance.
  \item \emph{Recall:} say the Japanese for firmly, then the Japanese for to follow an example, then say \emph{soroeru} again
\end{itemize}

\begin{tcolorbox}[breakable,colback=teal!4,colframe=teal!35!black,title={Before you move on}]
What does \ja{そろえる} mean? (\textquotedblleft{}To line up neatly; to put in order\textquotedblright{}.) Say it once more.
\end{tcolorbox}
\section[\texorpdfstring{\ja{あらためる} (aratameru)}{あらためる (aratameru)}]{\ja{あらためる} (aratameru) — to revise; to change for the better}

\label{lesson:JA-C330-aratameru}



\begin{quote}
Before the new one: say the Japanese for to follow an example, then the Japanese for to line up neatly.
\end{quote}

\subsection*{You'll want to know: \ja{あらためる}}
\textbf{\ja{あらためる}} — \emph{aratameru} — \textquotedblleft{}to revise; to change for the better\textquotedblright{}.

\begin{quote}
\textbf{\ja{きそくが} \ja{ふるいので、あらためた} \ja{ほうが} \ja{いいです。}} — \emph{kisoku ga furui node, aratameta hō ga ii desu} — \textquotedblleft{}The rules are out of date, so they should be revised.\textquotedblright{}
\end{quote}

Polite \textbf{\ja{あらためます}}, \emph{aratamemasu}; an ichidan verb. Past (polite): \textbf{\ja{かいしゃが} \ja{きそくを} \ja{あらためました。}}, \emph{kaisha ga kisoku o aratamemashita}.

\begin{cousinweb}
\textbf{\ja{あらためる}} (\emph{aratameru}) is a native Japanese word (wago), inherited rather than borrowed from Chinese, and its older origin is not settled. Its kanji was attached to it later as a kun'yomi, a native reading given to a Chinese character (JMdict, KANJIDIC2).
\end{cousinweb}

\begin{grammarlens}[title={What you've built}]
One more word for giving reasons and opinions.
\end{grammarlens}

\subsection*{Your turn}
\begin{itemize}
\raggedright
  \item \emph{Say it:} \emph{aratameru}
  \item \emph{Say it:} \emph{aratameru}, once more
  \item \emph{Say it:} \emph{kisoku ga furui node, aratameta hō ga ii desu} — The rules are out of date, so they should be revised.
  \item \emph{Recall:} say the Japanese for to follow an example, then the Japanese for to line up neatly, then say \emph{aratameru} again
\end{itemize}

\begin{tcolorbox}[breakable,colback=teal!4,colframe=teal!35!black,title={Before you move on}]
What does \ja{あらためる} mean? (\textquotedblleft{}To revise; to change for the better\textquotedblright{}.) Say it once more.
\end{tcolorbox}
